\documentclass[11pt,oneside]{article}

\pdfminorversion=7
\usepackage{QuizStyle}

\hypersetup{
    pdftitle={20261014 Quiz for MATH200 Differential Equations},
    pdfauthor={Jungwoo Kim},
    pdfsubject={Quiz for MATH200 Differential Equations},
    pdfcreator={LaTeX}
}

\begin{document}

\quiztitle{20261014 Quiz}

% Boyce Section 7.5, Problem 5(b), presented as a standalone problem.
\begin{problem}{7.5.5}
Find the general solution of the system
\[
    \mathbf{x}'=
    \begin{pmatrix}
        3 & -1\\
        6 & -2
    \end{pmatrix}\mathbf{x}.
\]
\end{problem}

\begin{solution}
Let
\[
    \mathbf{A}=\begin{pmatrix}3&-1\\6&-2\end{pmatrix}.
\]
The characteristic equation is
\[
    \det(\lambda\mathbf{I}_2-\mathbf{A})
    =\lambda^2-\lambda=\lambda(\lambda-1)=0.
\]
For the eigenvalue $\lambda_1=0$, solving $\mathbf{A}\mathbf{v}_1=\mathbf{0}$ gives
\[
    \mathbf{v}_1=\begin{pmatrix}1\\3\end{pmatrix}.
\]
For $\lambda_2=1$, solving $(\mathbf{A}-\mathbf{I}_2)\mathbf{v}_2=\mathbf{0}$ gives
\[
    \mathbf{v}_2=\begin{pmatrix}1\\2\end{pmatrix}.
\]
Since the eigenvalues are distinct, the corresponding eigenvectors are linearly independent. Therefore the general solution is
\[
    \boxed{
        \mathbf{x}(t)
        =c_1\begin{pmatrix}1\\3\end{pmatrix}
        +c_2e^t\begin{pmatrix}1\\2\end{pmatrix}
    }.
\]
The first term is constant because $e^{0t}=1$.
\end{solution}

\begin{problem}{7.5.11}
Solve the initial value problem
\[
    \mathbf{x}'=
    \begin{pmatrix}
        5&-1\\
        3&1
    \end{pmatrix}\mathbf{x},
    \qquad
    \mathbf{x}(0)=\begin{pmatrix}3\\-1\end{pmatrix}.
\]
Describe the behavior of the solution as $t\to\infty$.
\end{problem}

\begin{solution}
The characteristic equation of the coefficient matrix is
\[
    \det(\lambda\mathbf{I}_2-\mathbf{A})
    =\lambda^2-6\lambda+8=(\lambda-2)(\lambda-4)=0.
\]
The eigenvalues $\lambda_1=2$ and $\lambda_2=4$ have corresponding eigenvectors
\[
    \mathbf{v}_1=\begin{pmatrix}1\\3\end{pmatrix},
    \qquad
    \mathbf{v}_2=\begin{pmatrix}1\\1\end{pmatrix}.
\]
Thus
\[
    \mathbf{x}(t)
    =c_1e^{2t}\begin{pmatrix}1\\3\end{pmatrix}
     +c_2e^{4t}\begin{pmatrix}1\\1\end{pmatrix}.
\]
The initial condition gives
\[
    c_1+c_2=3,
    \qquad
    3c_1+c_2=-1,
\]
so $c_1=-2$ and $c_2=5$. Hence
\[
    \boxed{
        \mathbf{x}(t)
        =-2e^{2t}\begin{pmatrix}1\\3\end{pmatrix}
         +5e^{4t}\begin{pmatrix}1\\1\end{pmatrix}
    }.
\]
As $t\to\infty$, the $e^{4t}$ term dominates:
\[
    e^{-4t}\mathbf{x}(t)\longrightarrow
        5\begin{pmatrix}1\\1\end{pmatrix}.
\]
Thus both components tend to $+\infty$, and the trajectory moves away from the origin, asymptotically in the direction $(1,1)^{\mathsf T}$.
\end{solution}

\begin{problem}{7.5.12}
Solve the initial value problem
\[
    \mathbf{x}'=
    \begin{pmatrix}
        1&1&2\\
        0&2&2\\
        -1&1&3
    \end{pmatrix}\mathbf{x},
    \qquad
    \mathbf{x}(0)=\begin{pmatrix}2\\0\\3\end{pmatrix}.
\]
Describe the behavior of the solution as $t\to\infty$.
\end{problem}

\begin{solution}
For the coefficient matrix $\mathbf{A}$, the characteristic equation is
\[
    \det(\lambda\mathbf{I}_3-\mathbf{A})
    =(\lambda-1)(\lambda-2)(\lambda-3)=0.
\]
Solving $(\mathbf{A}-\lambda_i\mathbf{I}_3)\mathbf{v}_i=\mathbf{0}$ for $\lambda_1=1$, $\lambda_2=2$, and $\lambda_3=3$ gives
\[
    \mathbf{v}_1=\begin{pmatrix}0\\-2\\1\end{pmatrix},
    \qquad
    \mathbf{v}_2=\begin{pmatrix}1\\1\\0\end{pmatrix},
    \qquad
    \mathbf{v}_3=\begin{pmatrix}2\\2\\1\end{pmatrix}.
\]
Since the three eigenvalues are distinct, these eigenvectors form a basis of $\R^3$. The general solution is
\[
    \mathbf{x}(t)
    =c_1e^t\mathbf{v}_1+c_2e^{2t}\mathbf{v}_2+c_3e^{3t}\mathbf{v}_3.
\]
Applying the initial condition yields
\[
    c_2+2c_3=2,
    \qquad
    -2c_1+c_2+2c_3=0,
    \qquad
    c_1+c_3=3.
\]
Thus $c_1=1$, $c_2=-2$, and $c_3=2$, giving
\[
    \boxed{
    \begin{aligned}
        \mathbf{x}(t)
        ={}&e^t\begin{pmatrix}0\\-2\\1\end{pmatrix}
            -2e^{2t}\begin{pmatrix}1\\1\\0\end{pmatrix}\\
           &+2e^{3t}\begin{pmatrix}2\\2\\1\end{pmatrix}
    \end{aligned}
    }.
\]
The term with the largest eigenvalue dominates as $t\to\infty$:
\[
    e^{-3t}\mathbf{x}(t)\longrightarrow
    2\begin{pmatrix}2\\2\\1\end{pmatrix}.
\]
Therefore all three components tend to $+\infty$, and the trajectory moves away from the origin, asymptotically in the direction $(2,2,1)^{\mathsf T}$.
\end{solution}

\begin{problem}{7.6.4}
Consider the system
\[
    \mathbf{x}'=
    \begin{pmatrix}
        3&4\\
        -2&-1
    \end{pmatrix}\mathbf{x}.
\]
\begin{problemsubparts}
    \item Express the general solution of the given system in terms of real-valued functions.
    \item Describe the behavior of the solutions as $t\to\infty$.
\end{problemsubparts}
\end{problem}

\begin{solution}
\noindent\textup{(a)} The characteristic equation is
\[
    \det(\lambda\mathbf{I}_2-\mathbf{A})
    =\lambda^2-2\lambda+5=0,
\]
with roots $\lambda=1\pm2i$. For $\lambda=1+2i$, an eigenvector is
\[
    \mathbf{v}=\begin{pmatrix}-1-i\\1\end{pmatrix}
    =\begin{pmatrix}-1\\1\end{pmatrix}
      +i\begin{pmatrix}-1\\0\end{pmatrix}.
\]
By Euler's formula, the real and imaginary parts of the complex solution $e^{(1+2i)t}\mathbf{v}$ are
\[
\begin{aligned}
    \mathbf{x}^{(1)}(t)
       &=e^t\begin{pmatrix}
           \sin(2t)-\cos(2t)\\
           \cos(2t)
         \end{pmatrix},\\
    \mathbf{x}^{(2)}(t)
       &=e^t\begin{pmatrix}
           -\sin(2t)-\cos(2t)\\
           \sin(2t)
         \end{pmatrix}.
\end{aligned}
\]
Both are real-valued solutions, and
\[
    \det\bigl(\mathbf{x}^{(1)}(t)\ \mathbf{x}^{(2)}(t)\bigr)
    =e^{2t}\ne0.
\]
Hence they form a fundamental set and the real-valued general solution is
\[
    \boxed{
    \begin{aligned}
        \mathbf{x}(t)
        ={}&c_1e^t\begin{pmatrix}
            \sin(2t)-\cos(2t)\\
            \cos(2t)
        \end{pmatrix}\\
        &+c_2e^t\begin{pmatrix}
            -\sin(2t)-\cos(2t)\\
            \sin(2t)
        \end{pmatrix}
    \end{aligned}
    }.
\]

\medskip
\noindent\textup{(b)} The eigenvalues have positive real part, so every nonzero solution spirals away from the origin as $t\to\infty$. The rotation is clockwise because
\[
    x_1x_2'-x_2x_1'
    =-2\bigl((x_1+x_2)^2+x_2^2\bigr)<0
    \qquad (\mathbf{x}\ne\mathbf{0}).
\]
Hence the origin is an \(\boxed{\text{unstable clockwise spiral point}}\); the zero solution remains at the origin.
\end{solution}

\begin{problem}{7.7.2}
For the system
\[
    \mathbf{x}'=
    \begin{pmatrix}
        2&1\\
        -5&-2
    \end{pmatrix}\mathbf{x},
\]
\begin{problemsubparts}
    \item Find a fundamental matrix.
    \item Find the fundamental matrix $\boldsymbol{\Phi}(t)$ satisfying $\boldsymbol{\Phi}(0)=\mathbf{I}_2$.
\end{problemsubparts}
\end{problem}

\begin{solution}
\noindent\textup{(a)} The characteristic equation is
\[
    \det(\lambda\mathbf{I}_2-\mathbf{A})=\lambda^2+1=0,
\]
so $\lambda=\pm i$. An eigenvector for $\lambda=i$ is
\[
    \mathbf{v}=\begin{pmatrix}1\\-2+i\end{pmatrix}.
\]
Taking the real and imaginary parts of $e^{it}\mathbf{v}$ gives
\[
    \mathbf{x}^{(1)}(t)=
    \begin{pmatrix}\cos t\\-2\cos t-\sin t\end{pmatrix},
    \qquad
    \mathbf{x}^{(2)}(t)=
    \begin{pmatrix}\sin t\\\cos t-2\sin t\end{pmatrix}.
\]
Placing these independent solutions in the columns gives
\[
    \boxed{
        \boldsymbol{\Psi}(t)=
        \begin{pmatrix}
            \cos t & \sin t\\
            -2\cos t-\sin t & \cos t-2\sin t
        \end{pmatrix}
    }.
\]
Indeed, $\det\boldsymbol{\Psi}(t)=1\ne0$, so $\boldsymbol{\Psi}$ is a fundamental matrix.

\medskip
\noindent\textup{(b)} To impose the normalization at $t=0$, use
\[
    \boldsymbol{\Phi}(t)
    =\boldsymbol{\Psi}(t)\boldsymbol{\Psi}(0)^{-1}.
\]
Since
\[
    \boldsymbol{\Psi}(0)
    =\begin{pmatrix}1&0\\-2&1\end{pmatrix},
    \qquad
    \boldsymbol{\Psi}(0)^{-1}
    =\begin{pmatrix}1&0\\2&1\end{pmatrix},
\]
we obtain
\[
    \boxed{
        \boldsymbol{\Phi}(t)=
        \begin{pmatrix}
            \cos t+2\sin t & \sin t\\
            -5\sin t & \cos t-2\sin t
        \end{pmatrix}
    }.
\]
Direct substitution confirms
\[
    \boldsymbol{\Phi}'(t)=\mathbf{A}\boldsymbol{\Phi}(t),
    \qquad
    \boldsymbol{\Phi}(0)=\mathbf{I}_2,
\]
as required.
\end{solution}

\end{document}
